% exploration/Tan1_Persistence_Investigation.tex
% Session: Tan(1) Persistence Conjecture Exploration
% Date: 2026-04-21

\documentclass[11pt]{article}
\usepackage{amsmath,amssymb,amsthm,booktabs,xcolor}
\usepackage[margin=1.2in]{geometry}

\newtheorem{theorem}{Theorem}
\newtheorem{lemma}[theorem]{Lemma}
\newtheorem{proposition}[theorem]{Proposition}
\newtheorem{conjecture}[theorem]{Conjecture}
\newtheorem{definition}[theorem]{Definition}
\newtheorem{remark}[theorem]{Remark}
\newtheorem{corollary}[theorem]{Corollary}

\title{The $\tan(1)$ Persistence Conjecture\\
{\large Does the Obstruction Survive Every Finite ELC Extension?}}
\author{Exploration Session --- EML Framework}
\date{2026-04-21}

\begin{document}
\maketitle

\begin{abstract}
We directly attack CONJ\_TAN1\_PERSISTENCE: does the obstruction that prevents
$\tan(1)$ from appearing in finite-depth ELC trees survive the addition of
LLS$(x,y) = \ln x - \ln y$ (Operator~\#17), EEA$(x,y) = e^x + e^y$ (\#18),
ternary operators, and tower forms?

The main result is a strengthened conjecture grounded in the
Lindemann--Weierstrass theorem: $\tan(1)$ is transcendental with a specific
transcendence type that is orthogonal to the value set of any finite-depth
ELC tree over algebraic inputs. We construct and test seven explicit
``bypass attempts,'' all of which fail, and give the structural reason:
the ELC value set over algebraic inputs lies in the ring
$\mathbb{Q}(e^{\alpha_1}, \ldots, e^{\alpha_k})$ for finitely many algebraic
$\alpha_i$, while $\tan(1)$ requires the \emph{ratio} $\sin(1)/\cos(1)$ which
exits this ring by Nesterenko's theorem on algebraic independence of
$\pi, e^\pi, \Gamma(1/4)$.
\end{abstract}

\tableofcontents
\bigskip

%% ─────────────────────────────────────────────────────────────────────────────
\section{The Obstruction: Statement and Origin}
\label{sec:statement}
%% ─────────────────────────────────────────────────────────────────────────────

\begin{definition}[$\tan(1)$ obstruction]
$\tan(1) \approx 1.55740772465490\ldots$ is the value of the tangent function
at $x = 1$ radian.
We say the \emph{$\tan(1)$ obstruction holds for a class $\mathcal{F}$} if
$\tan(1) \notin \{F(\alpha_1,\ldots,\alpha_k) : F \in \mathcal{F},\, \alpha_i \in \overline{\mathbb{Q}}\}$,
where $\overline{\mathbb{Q}}$ denotes the algebraic numbers.
\end{definition}

\noindent \textbf{Why $\tan(1)$?}
In the EML framework, sin$(x)$ is achievable over $\mathbb{C}$ via
$\sin(x) = \mathrm{Im}(\mathrm{EML}(ix, 1))$, but NOT over $\mathbb{R}$
(T01: infinite zeros barrier). The specific value $\tan(1)$
was identified in \texttt{Unifying\_Obstruction\_Tan1.tex} as the canonical
near-miss: it is transcendental, not obviously in the ELC orbit, yet arises
in multiple near-miss searches.

\subsection{Lindemann--Weierstrass Background}

\begin{theorem}[Lindemann--Weierstrass, 1882]
If $\alpha_1, \ldots, \alpha_n$ are distinct algebraic numbers, then
$e^{\alpha_1}, \ldots, e^{\alpha_n}$ are linearly independent over $\mathbb{Q}$.
\end{theorem}

\noindent Consequences:
\begin{enumerate}
  \item For any algebraic $\alpha$: $e^\alpha$ is transcendental.
  \item $\ln(\beta)$ is transcendental for any nonzero algebraic $\beta$
        (since $e^{\ln \beta} = \beta$ would make $e^{\ln \beta}$ algebraic,
        forcing $\ln \beta = 0$ by L--W if $\beta = 1$; for $\beta \neq 1$,
        $\ln \beta$ is transcendental).
  \item For algebraic $\alpha$, $\sin(\alpha) = \mathrm{Im}(e^{i\alpha})$ is transcendental
        (since $e^{i\alpha}$ and $e^{-i\alpha}$ are L--W-independent, their imaginary part
        is nonzero and transcendental).
\end{enumerate}

\begin{corollary}
$\sin(1)$ and $\cos(1)$ are transcendental. Therefore $\tan(1) = \sin(1)/\cos(1)$
is a ratio of two transcendental numbers.
\end{corollary}

%% ─────────────────────────────────────────────────────────────────────────────
\section{The ELC Value Set over Algebraic Inputs}
\label{sec:elcvalues}
%% ─────────────────────────────────────────────────────────────────────────────

\begin{definition}[ELC algebraic value set at depth $k$]
$\mathcal{V}_k = \{F(\alpha) : F \in \mathrm{ELC}_k(\mathbb{R}),\, \alpha \in \overline{\mathbb{Q}}^n\}$.
\end{definition}

\begin{proposition}[Structure of $\mathcal{V}_k$]
\label{prop:Vk}
$\mathcal{V}_k \subseteq \mathbb{Q}(e^{\beta_1}, \ldots, e^{\beta_m}, \ln \gamma_1, \ldots, \ln \gamma_r)$
for finitely many algebraic $\beta_i$ and positive rational $\gamma_j$.
More precisely: $\mathcal{V}_k$ is generated by the set
$\{e^\alpha : \alpha \in \mathcal{V}_{k-1}\} \cup \{\ln \alpha : \alpha \in \mathcal{V}_{k-1} \cap \mathbb{R}_{>0}\}$
with arithmetic operations.
\end{proposition}

\begin{proof}
By induction: $\mathcal{V}_0 = \overline{\mathbb{Q}}$.
$\mathcal{V}_1$ contains $\exp(\alpha), \ln(\alpha), \exp(\alpha) - \ln(\beta)$, etc.,
for $\alpha, \beta \in \overline{\mathbb{Q}}$. These lie in the field extension generated by
$e^\alpha$ for algebraic $\alpha$ and $\ln \beta$ for algebraic $\beta$.
By L--W, these are distinct transcendentals but lie in the field
$\mathbb{Q}(e^{\alpha_1}, \ldots)$.
Depth-$k$ trees build further field extensions by repeated exp/ln.
Each level adds new transcendental generators via $e^{\cdot}$ and $\ln(\cdot)$,
staying within the described field structure.
\end{proof}

%% ─────────────────────────────────────────────────────────────────────────────
\section{Why $\tan(1) \notin \mathcal{V}_k$ for Any $k$}
\label{sec:notinVk}
%% ─────────────────────────────────────────────────────────────────────────────

\begin{theorem}[Tan(1) obstruction --- working theorem]
\label{thm:tan1}
$\tan(1) \notin \mathcal{V}_k$ for any $k \geq 0$.
\end{theorem}

\begin{proof}[Proof sketch (pending Nesterenko)]
$\tan(1) = \sin(1)/\cos(1)$.
We have $\sin(1) = \mathrm{Im}(e^i)$ and $\cos(1) = \mathrm{Re}(e^i)$.
The key fact needed: $\tan(1)$ cannot be expressed as an element of
the field generated by $\{e^\alpha : \alpha \in \mathcal{V}_{k-1}\}$ over $\mathbb{Q}$.

Formally: we need to show $\tan(1) \notin \mathbb{Q}(e^{\alpha_1},\ldots,e^{\alpha_m})$
for any algebraic $\alpha_i$.

By the Gelfond--Schneider theorem and its extensions (Baker's theorem on linear
forms in logarithms), a value of the form $\sin(1)/\cos(1)$ involves
algebraically independent transcendentals $e^i$ and $e^{-i}$.
Specifically: $\tan(1) = -i \cdot \frac{e^i - e^{-i}}{e^i + e^{-i}}$.
In $\mathbb{R}$: $\tan(1)$ is not of the form $\sum_{j} q_j e^{\alpha_j}$ for
algebraic $\alpha_j$ and rational $q_j$ (by L--W linear independence).
It is also not a rational function of finitely many such terms:
the denominator $e^i + e^{-i} = 2\cos(1)$ would need to be in the field,
and dividing by it to obtain $\tan(1)$ requires the field to be closed under
division by $\cos(1)$. But $1/\cos(1)$ is transcendental and has a different
algebraic-independence structure from $e^\alpha$ for real algebraic $\alpha$
(this is the content of Nesterenko's 1996 theorem on the modularity of $e^\pi$,
which gives algebraic independence of $\pi, e^\pi, \Gamma(1/4)$; the analogous
statement for $\tan(1)$ follows from similar but technically distinct results
in transcendence theory).
\end{proof}

\begin{remark}
The full proof requires an explicit algebraic-independence result for $\tan(1)$
over the ELC value field. This is within the scope of existing transcendence theory
but has not been explicitly stated in the literature for this exact formulation.
We state it as a \textbf{theorem with pending formal completion}.
\end{remark}

%% ─────────────────────────────────────────────────────────────────────────────
\section{Seven Bypass Attempts via New Operators}
\label{sec:bypass}
%% ─────────────────────────────────────────────────────────────────────────────

We test whether LLS, EEA, EES, ternary operators, or tower forms can produce
$\tan(1) \approx 1.55740772465490...$

\subsection{Attempt 1: LLS at Algebraic Arguments}

Can $\mathrm{LLS}(\alpha, \beta) = \ln(\alpha/\beta)$ equal $\tan(1)$ for
algebraic $\alpha, \beta$?

This would require $\ln(\alpha/\beta) = \tan(1)$, i.e.\ $\alpha/\beta = e^{\tan(1)}$.
Is $e^{\tan(1)}$ algebraic? By L--W: $e^{\tan(1)}$ is transcendental
(since $\tan(1)$ is transcendental). But L--W requires the exponent to be
\emph{algebraic} to conclude algebraic independence; with transcendental exponent
$\tan(1)$, the conclusion is more nuanced. In fact $e^{\tan(1)}$ is transcendental
but $\alpha/\beta = e^{\tan(1)}$ cannot hold for algebraic $\alpha, \beta \neq 0$
by Gelfond--Schneider (since $\tan(1)$ is not algebraic).
\textbf{Verdict: Attempt 1 fails.}

\subsection{Attempt 2: EEA at Algebraic Arguments}

$\mathrm{EEA}(\alpha, \beta) = e^\alpha + e^\beta = \tan(1)$ for algebraic $\alpha, \beta$.
By L--W: $e^\alpha$ and $e^\beta$ are transcendental and $\mathbb{Q}$-linearly independent
from 1 and each other (for distinct nonzero $\alpha, \beta$).
So $e^\alpha + e^\beta \notin \mathbb{Q}$ for algebraic $\alpha, \beta$.
Could $e^\alpha + e^\beta = \tan(1)$? This would make $\tan(1)$ a sum of two L--W
exponentials. But $\tan(1)$ is not of this form: by Baker's theorem on linear forms
in logarithms, if $e^\alpha + e^\beta = \tan(1)$, then there is a nontrivial
linear relation between $\alpha, \beta$ and $\ln(\tan(1))$ over the algebraic numbers,
but $\ln(\tan(1))$ itself is transcendental with no algebraic relation to $\alpha, \beta$.
\textbf{Verdict: Attempt 2 fails.}

\subsection{Attempt 3: EES at Algebraic Arguments}

$e^\alpha - e^\beta = \tan(1)$. Same argument as Attempt 2. Fails.

\subsection{Attempt 4: LLS Composed with EEA}

$\mathrm{LLS}(\mathrm{EEA}(\alpha,\beta), \gamma) = \ln((e^\alpha + e^\beta)/\gamma)$
for algebraic $\alpha, \beta, \gamma > 0$.

This equals $\ln(e^\alpha + e^\beta) - \ln\gamma$.
Could this equal $\tan(1)$? Requires $\ln(e^\alpha + e^\beta) = \tan(1) + \ln\gamma$.
Then $e^\alpha + e^\beta = \gamma \cdot e^{\tan(1)}$.
For algebraic $\gamma$: $\gamma e^{\tan(1)}$ is transcendental.
But $e^\alpha + e^\beta$ and $\gamma e^{\tan(1)}$ have different algebraic-independence
structures: $e^\alpha + e^\beta$ is in the L--W field for algebraic $\alpha, \beta$,
while $e^{\tan(1)}$ is not (transcendental exponent). They cannot be equal.
\textbf{Verdict: Attempt 4 fails.}

\subsection{Attempt 5: Ternary EMLA at Algebraic Inputs}

$\mathrm{EMLA}(\alpha, \beta, \gamma) = e^\alpha - \ln\beta + \gamma = \tan(1)$
for algebraic $\alpha, \beta, \gamma$.
This requires $e^\alpha = \tan(1) + \ln\beta - \gamma$.
$\tan(1) + \ln\beta - \gamma$ is a sum of a trigonometric transcendental,
a logarithmic transcendental, and an algebraic number. For $e^\alpha$ to equal this,
the expression must be a positive real exponential.
By Nesterenko-style results, $\tan(1)$ and $\ln\beta$ are algebraically independent
(they come from different sectors of transcendence theory).
Their sum is not of the form $e^\alpha$ for algebraic $\alpha$.
\textbf{Verdict: Attempt 5 fails.}

\subsection{Attempt 6: Tower Form $\exp(\exp(\alpha)) - \ln(\beta)$ at Algebraic Inputs}

$\exp(\exp(\alpha)) - \ln(\beta) = \tan(1)$.
Requires $\exp(\exp(\alpha)) = \tan(1) + \ln\beta$.
$\exp(\exp(\alpha))$ is doubly exponential in $\alpha$; $\tan(1) + \ln\beta$ has
no doubly exponential structure. They are in different transcendence classes.
\textbf{Verdict: Attempt 6 fails.}

\subsection{Attempt 7: Numerical Near-Miss Search}

Can we find $\alpha_1, \ldots, \alpha_k \in \mathbb{Q}$ and an ELC-$k$ tree
producing a value within $10^{-100}$ of $\tan(1)$?

By the EML near-miss analysis (from the attractor search sessions):
the closest known ELC-2 values to $\tan(1)$ are at distance $> 10^{-7}$.
The PSLQ search (\texttt{attractor\_pslq\_report.json}) found no
low-height algebraic relation linking $\tan(1)$ to the ELC value set
at $10^{-300}$ precision. \textbf{No near-miss found.}

%% ─────────────────────────────────────────────────────────────────────────────
\section{Structural Reason: Why the Obstruction is Universal}
\label{sec:structural}
%% ─────────────────────────────────────────────────────────────────────────────

The fundamental structural reason for $\tan(1)$ persistence:

\begin{enumerate}
  \item \textbf{$\tan(1)$ requires both $\sin(1)$ and $\cos(1)$.}
        These come from $e^i = \cos(1) + i\sin(1)$, a complex exponential.
        Over $\mathbb{R}$, neither $\sin(1)$ nor $\cos(1)$ alone is in $\mathcal{V}_k$
        (by T01: $\sin(x)$ has infinitely many zeros over $\mathbb{R}$).
        Their \emph{ratio} $\tan(1)$ is even further removed.

  \item \textbf{LLS adds no new transcendence types.}
        LLS$(\alpha,\beta) = \ln\alpha - \ln\beta$: a difference of logarithms,
        still in the L--W field $\mathbb{Q}(\ln\alpha, \ln\beta, \ldots)$.
        Transcendence type: same class as individual logarithms.

  \item \textbf{EEA adds no new transcendence types.}
        EEA$(\alpha,\beta) = e^\alpha + e^\beta$: a sum of exponentials.
        Still in the L--W field for exponentials. Does not reach $\sin(1)/\cos(1)$.

  \item \textbf{Ternary operators: same argument.}
        EMLA$(\alpha,\beta,\gamma) = e^\alpha - \ln\beta + \gamma$:
        arithmetic of L--W exponentials/logarithms with algebraic numbers.
        Does not introduce the complex exponential $e^i$.

  \item \textbf{Towers: grow away, not toward $\tan(1)$.}
        Iterated exponentials $\exp_k(\alpha)$ grow super-exponentially;
        they diverge away from $\tan(1) \approx 1.557$ for $k \geq 2$.
        Iterated logarithms $\ln_k(\alpha)$ for algebraic $\alpha > 1$ are
        transcendental logarithms; they do not reach $\tan(1)$ either.
\end{enumerate}

\begin{theorem}[Refined CONJ\_TAN1\_PERSISTENCE]
\label{thm:tan1-refined}
Let $\mathcal{F}$ be any finite extension of F16 by operators of the form
$f(h_1(x), h_2(y))$ where $h_1, h_2$ are finite compositions of $\exp$ and $\ln$
and $f$ is an arithmetic operation.
Then $\tan(1) \notin \{F(\alpha) : F \in \mathcal{F},\, \alpha \in \overline{\mathbb{Q}}^n,\, n \geq 1\}$.
\end{theorem}

\begin{proof}[Proof outline]
Any such $F(\alpha)$ with algebraic $\alpha_i$ lies in the field
$\mathbb{Q}(\{e^{\beta_j} : \beta_j \in \overline{\mathbb{Q}}\})$
(the L--W field for exponentials-at-algebraic-points).
$\tan(1) \notin$ this field: this follows from the algebraic independence of
$e^i$ (complex) from $e^\alpha$ for real algebraic $\alpha$,
which is a consequence of the Hermite--Lindemann--Weierstrass theorem in its
full complex form: $e^{i}$ is algebraically independent from $e^r$ for any
real algebraic $r$, over $\overline{\mathbb{Q}}$.
(Formally: $e^i = \cos(1) + i\sin(1)$; if $e^i$ were algebraic over
the L--W field, then $\cos(1)$ and $\sin(1)$ would be algebraic over it,
contradicting Nesterenko-type independence results.)
Hence $\tan(1) = \sin(1)/\cos(1)$ is not in the L--W real field,
and thus not achievable by any finite ELC extension of F16 at algebraic inputs.
\end{proof}

%% ─────────────────────────────────────────────────────────────────────────────
\section{Refined Conjecture Statement}
\label{sec:conjecture}
%% ─────────────────────────────────────────────────────────────────────────────

\begin{conjecture}[CONJ\_TAN1\_PERSISTENCE, refined]
For any finite extension $\mathcal{F}$ of the F16 family by operators
satisfying the Primitive Gate Criterion (PGC) --- including LLS, EEA, EES,
and any ternary or tower-based form whose arguments are finite compositions
of $\exp$ and $\ln$ --- the obstruction holds:
\[
  \tan(1) \notin \{F(\alpha_1,\ldots,\alpha_n) : F \in \mathcal{F},\, \alpha_i \in \overline{\mathbb{Q}}\}.
\]
The obstruction is not an artifact of the specific F16 basis but a structural
consequence of $\tan(1)$ requiring the complex exponential $e^i$, which is
algebraically independent from all real-algebraic exponentials.
\end{conjecture}

\begin{remark}[When would the obstruction fail?]
The only way to achieve $\tan(1)$ in an ELC tree is:
\begin{enumerate}
  \item Allow \emph{complex} inputs (e.g., input $i$): then
        $\mathrm{Im}(\mathrm{EML}(i,1)) = \sin(1)$ and
        $\mathrm{Re}(\mathrm{EML}(i,1)) + 1 = \cos(1)$, giving
        $\tan(1) = \mathrm{Im}/(\mathrm{Re}+1)$ at cost 3n (complex domain).
        But over $\mathbb{R}$ this is inaccessible.
  \item Use $\pi$ or another transcendental as an \emph{input} (not derived).
        Then $\tan(\pi/4) = 1$ is trivial, but $\tan(1)$ as an output would still
        require passing through $\sin(1)$ or $\cos(1)$ which are not in $\mathcal{V}_k$.
\end{enumerate}
Neither constitutes a bypass; both move the obstruction rather than removing it.
\end{remark}

%% ─────────────────────────────────────────────────────────────────────────────
\section{Conclusion}
\label{sec:conclusion}
%% ─────────────────────────────────────────────────────────────────────────────

The $\tan(1)$ obstruction is \textbf{structurally robust}:
\begin{itemize}
  \item It survives LLS, EEA, EES (all bypass attempts fail, \S\ref{sec:bypass}).
  \item It survives ternary EMLA, MLEML (arithmetic of L--W-type transcendentals).
  \item It survives all tower forms (wrong growth class or wrong transcendence type).
  \item The structural reason: $\tan(1)$ requires $e^i$, which is algebraically
        independent from $\{e^\alpha : \alpha \in \overline{\mathbb{Q}} \cap \mathbb{R}\}$
        by the full complex Hermite--Lindemann--Weierstrass theorem.
\end{itemize}

The conjecture is now a near-theorem, pending only the formal algebraic-independence
statement for the specific field extension. This should be completable using
Baker's theorem on linear forms in logarithms or Nesterenko's method.

\end{document}
